\section{Factorial divisors and weighted $p$-adic root counts}
\label{sec:div-discs}
The appropriate local quantity counts roots in discs centred at zero
and integrates over their radii. Counting all integral roots once at
each prime would give a divergent expression even for one integer root.

Let $F\in\mathbb Z[X]$ be nonzero with $F(0)\ne0$, let $c(F)>0$
be the greatest common divisor of its coefficients, and normalize $v_p(p)=1$.
Roots in $\overline{\mathbb Q}_p$ are counted with multiplicity.
For $t>0$ put
\[
 n_{F,p}(t)=\#\{\alpha:F(\alpha)=0,\ v_p(\alpha)\ge t\},
 \qquad I_p(F)=\int_0^\infty n_{F,p}(t)\,dt.
\]
The endpoints of the finitely many intervals do not affect the integral.

\begin{proposition}\label{prop:disc-product-formula}
For every such $F$,
\[
 I_p(F)=v_p(F(0))-v_p(c(F)),\qquad
 \sum_p I_p(F)\log p=\log\frac{|F(0)|}{c(F)}.
\]
If $D_N(X)=\prod_{j=1}^N(X+j)$ divides $F$ in $\mathbb Z[X]$, then
\[
 \log N!\ \le\ \sum_p I_p(F)\log p
 \ \le\ \log M(F)-\log c(F),
\]
where $M(F)=|a|\prod_\alpha\max(1,|\alpha|)$ is its complex
Mahler measure and $a$ its leading coefficient.
For the divisor itself,
\[
 I_p(D_N)=\sum_{k\ge1}\left\lfloor N/p^k\right\rfloor=v_p(N!),
 \qquad M(D_N)=N!.
\]
\end{proposition}
\begin{proof}
Write $F=a\prod_i(X-\alpha_i)$ over an algebraic closure of
$\mathbb Q_p$. The Gauss norm is multiplicative: after scaling two
polynomials to have integral coefficients with at least one unit,
their nonzero reductions multiply to a nonzero polynomial over the
residue field. Applying this to the linear factors yields
\[
 \min_jv_p([X^j]F)=v_p(a)+\sum_i\min(0,v_p(\alpha_i)).
\]
The left side is $v_p(c(F))$, whereas
$v_p(F(0))=v_p(a)+\sum_i v_p(\alpha_i)$.
Subtracting proves the first identity, since integration of the
indicator of $0<t\le v_p(\alpha_i)$ gives its positive part.
The rational product formula gives the second identity; only finitely
many primes contribute because both $F(0)$ and $c(F)$ are integers.

Write $F=D_NH$. Since $D_N$ is primitive, Gauss's lemma gives
$c(F)=c(H)$. Thus
$|F(0)|/c(F)=N!\,|H(0)|/c(H)\ge N!$.
The complex root product gives $|F(0)|\le M(F)$ directly.
Finally the roots of $D_N$ are the nonzero integers $-j$.
Their valuations are integers, and counting the multiples of $p^k$
gives the displayed floor sum. Their archimedean positive parts
give $M(D_N)=\prod_{j=1}^Nj$.
\end{proof}

\paragraph{The Newton-polygon convention.}
The root valuations used below follow from the same Gauss-norm
argument. If $F=\sum_j a_jX^j=a_d\prod_i(X-\alpha_i)$, then
for every real $t$,
\[
 \min_j\{v_p(a_j)+jt\}
 =v_p(a_d)+\sum_i\min\{t,v_p(\alpha_i)\}.
\]
For rational $t$, choose a scalar $\beta\in\overline{\mathbb Q}_p$
with $v_p(\beta)=t$ and apply Gauss multiplicativity to $F(\beta X)$.
Both sides are continuous piecewise-linear functions, so the identity
extends to real $t$. A root of valuation $v$ changes the right-hand
slope by one as $t$ crosses $v$. On the left this transition occurs
when $t$ is the negative of a slope of the lower convex hull of
the coefficient points $(j,v_p(a_j))$, with zero coefficients omitted.
Consequently the negatives of those polygon slopes are the root
valuations, and the horizontal lengths are their multiplicities.
This convention includes fractional valuations and justifies the
disc counts computed below.

If $F=\sum_{i=1}^{K}\epsilon_i f_i^M$ with $|\epsilon_i|\le1$,
and each integral $f_i$ has at most $T$ monomials of coefficient
absolute value at most $H$, then $\|F\|_1\le K(TH)^M$.
The elementary circle-mean formula for each linear
factor gives $\log M(F)=\int_0^1\log|F(e^{2\pi it})|\,dt$,
hence $M(F)\le\|F\|_1$. Therefore
\[
 \log N!\le\log K+M\log(TH)-\log c(F).
\]
This retains a coefficient-height parameter. The local identity
does not eliminate it.

\subsection{Six distinct integer roots}
Consider the sparse integer polynomials
\[
 u=360-342X+42X^3,\qquad v=119X^2-60X^3+X^4.
\]
Direct multiplication gives the complete identity
\[
 u^2-v^2=-\prod_{j=1}^6(X-j)(X^2-99X-180).
\]
In particular the polynomial is nonzero and primitive. Replacing
$X$ by $-X$ gives a rising-factorial divisor, without changing any
root valuation or the Mahler measure.
The six integer roots contribute $v_p(6!)$; the quadratic factor
has positive root valuations $(2)$ at $p=2$, $(1,1)$ at $p=3$,
and $(1)$ at $p=5$. These follow from the lower Newton polygons
of the coefficient triples $(-180,-99,1)$.
Consequently the nonzero integer-radius disc counts are
\[
\begin{array}{c|c|c|c}
 p&n_{F,p}(1)&n_{F,p}(2)&I_p(F)\\\hline
 2&4&2&6\\
 3&4&0&4\\
 5&2&0&2
\end{array}
\]
and all higher integer-radius counts vanish. Thus
\[
 \sum_p I_p(F)\log p
 =6\log2+4\log3+2\log5=\log129600=\log(6!\cdot180).
\]

\subsection{A root of multiplicity sixteen}
Now consider the sparse integer polynomials
\begin{align*}
 u={}&-57+912X^5-969X^6+969X^{16}-912X^{17}+57X^{22},\\
 v={}&272X-374X^2+646X^{10}-646X^{12}+374X^{20}-272X^{21}.
\end{align*}
Define two reciprocal polynomials by their ascending coefficient lists:
\begin{align*}
 A:\;&(57,298,846,1618,2331,2620,2331,1618,846,298,57),\\
 B:\;&(57,500,1341,2864,5036,7184,9287,10684,11366,\\
     &\hspace{20mm}10684,9287,7184,5036,2864,1341,500,57).
\end{align*}
The finite multiplication identities
\[
 u+v=(X-1)^{11}(X+1)A(X),\qquad
 u-v=(X-1)^5(X+1)B(X)
\]
give
\[
 F=u^2-v^2=(X-1)^{16}(X+1)^2A(X)B(X).
\]
The first nonzero coefficient after $X=1+Y$ is
\[
 [Y^{16}]F(1+Y)=4A(1)B(1)=4406856960
 =2^8\cdot3\cdot5\cdot11\cdot17^2\cdot19^2.
\]
It is nonzero, so the order is exactly sixteen.
This is a repeated-root example; it does not assert a divisor of
sixteen distinct consecutive roots and does not give a $16!$ bound.

To make a finite local count at the repeated root, remove that root
and put $H(Y)=F(1+Y)/Y^{16}$. It is primitive: the constant and $X^2$ coefficients of $F$ are $3249$ and $-73984$, which are coprime; integral translation preserves content, as does removal of a monomial factor.
Its positive root valuations, listed as (valuation, multiplicity), are
\[
\begin{array}{c|l|r}
 p&\text{positive valuations of roots of }H&I_p(H)\\\hline
 2&(1,4),\ (1/4,16)&8\\
 3&(1/6,6)&1\\
 5&(1/4,4)&1\\
 11&(1/10,10)&1\\
 17&(1/6,6),\ (1/12,12)&2\\
 19&(1/8,8),\ (1/14,14)&2
\end{array}
\]
For example, at two the lower hull has vertices
$(0,8),(4,4),(20,0),(28,0)$; its first two slopes give the first row.
At $3,5,11,17,19$ the negative-slope vertices are respectively
\[
\begin{gathered}
 (0,1),(6,0);\quad(0,1),(4,0);\quad(0,1),(10,0);\\
 (0,2),(6,1),(18,0);\quad(0,2),(8,1),(22,0).
\end{gathered}
\]
The remaining vertices have nonnegative slopes. These coefficient
and hull computations follow by expanding the two displayed factor
identities; they involve integer arithmetic and rational comparisons.
It follows that
\[
 \sum_p I_p(H)\log p=\log4406856960.
\]
The fractional valuations explain why a sum only over discs
$p^k\mathbb Z_p$ would lose information here. The integration is over
all positive real valuation thresholds in $\overline{\mathbb Q}_p$.
